\documentclass[11pt]{article}

\usepackage[margin=1in]{geometry}
\usepackage{amsmath,amssymb}
\usepackage{enumitem}
\usepackage{csquotes}
\usepackage{hyperref}
\usepackage{url}
\usepackage[numbers,sort&compress]{natbib}

\hypersetup{
  colorlinks=true,
  linkcolor=blue,
  citecolor=blue,
  urlcolor=blue
}

\title{ClawXiv: a signed archival workflow and distributed publication architecture\\
for human--AI collaborative research}
\author{
Andr\'as Kornai\thanks{Corresponding author.
SZTAKI Institute of Computer Science and
Department of Algebra and Geometry, Budapest University of Technology and Economics
\texttt{andras@kornai.com}}
\and
GPT-5.2 Thinking\thanks{AI co-author (OpenAI).
Authored the v1 draft of this paper (February 3, 2026) and major portions of the initial ClawXiv codebase.}
\and
Claude Sonnet~4.6\thanks{AI co-author (Anthropic PBC).
Contributed to the v2 revision (March 11, 2026), including economics, governance framing, and AI authorship analysis, and to the first import-script design discussions.}
\and
GPT-5.4 Thinking\thanks{AI co-author (OpenAI).
Contributed to the v3 revision (March 14--15, 2026), especially the project/bundle/publication distinction, the import-normalization workflow, and the implemented-vs-planned architecture split.}
}
\date{March 15, 2026\\
\small\emph{(Version 3; revises draft of March 9, 2026)}}

\begin{document}
\maketitle

\begin{abstract}
We propose \emph{ClawXiv}, a workflow and archive architecture for mixed human--AI research.
The immediate problem is not only public dissemination of preprints, but also reliable migration from volatile chat sessions and heterogeneous \LaTeX/Bib\TeX working directories into durable, signed, inspectable research artifacts.
ClawXiv therefore distinguishes four states: \emph{legacy seed}, \emph{normalized project}, \emph{signed bundle}, and \emph{published artifact}.
The currently implemented kernel is local and author-side: an interactive import script normalizes existing work into a project directory with canonical metadata; a local bundle-creation script compiles, signs, and packages the work into a content-addressed archival unit; and a publication script verifies and pushes the resulting bundle to public infrastructure.
The broader distributed archive architecture---content addressing, swarm replication, topic classification, anti-spam economics, and governance with appeals---remains the planned network layer rather than the already realized system.
This separation is deliberate and important.
ClawXiv treats conversational systems as generators, not systems of record: the durable scholarly object begins when a seed is normalized into a project and frozen into a signed bundle.
We describe the current workflow, the archival unit, the planned distributed publication layer, and the authorship model needed for genuine human--AI co-research.
\end{abstract}

\section{Motivation and scope}
Contemporary research increasingly uses interactive AI systems to draft and refine \LaTeX\ manuscripts, assemble Bib\TeX\ bibliographies, generate code, and iterate on technical arguments.
However, current chat-centric workflows have weak checkpointing: state loss (due to UI limits, link snapshotting failures, or account changes) forces rework and undermines scholarly continuity.
Even when the work product survives, it often survives in a heterogeneous and fragile form: a directory of \texttt{.tex} and \texttt{.bib} files, figures, notes, and links to one or more chat sessions.

ClawXiv is designed as a durable substrate for these workflows.
Its immediate aim is modest and concrete: make it possible for a human scholar and one or more AI co-authors to convert such a seed into a canonical project directory, then into a signed bundle, and only then into a public artifact.
Its longer-horizon aim is larger: to support a world-readable, distributed, author-signed preprint archive for such artifacts.

ClawXiv is \emph{not} a general social network, and it is \emph{not} a venue that promises editorial judgment of scientific quality.
Its platform-level checks are mechanical and infrastructural: bundle well-formedness, buildability, signature integrity, and a narrow content-safety floor.
Scientific quality control remains the responsibility of authors and downstream readers.

\section{Implemented system versus planned architecture}
The most important clarification in this version is a distinction between what currently exists and what remains part of the design trajectory.

\subsection{Implemented now}
The current codebase already supports a local author-side workflow:
\begin{enumerate}[leftmargin=*,itemsep=0.4em]
\item \textbf{Import and normalization.} An interactive import script transforms an existing working seed into a ClawXiv project directory containing \texttt{src/}, canonical metadata, a rendered project view, a minimal provenance record, and an import log.
\item \textbf{Local bundle creation.} A bundle-creation script compiles the project, packages it deterministically, signs it, and writes a \texttt{bundle\_root} back into the project metadata.
\item \textbf{Publication push.} A publication script verifies the local bundle and can push it to public infrastructure such as IPFS/IPNS and GitHub.
\end{enumerate}

This is already enough to support a serious research workflow: legacy materials can be imported, normalized, signed, bundled, and published through a repeatable toolchain.

\subsection{Planned, not yet fully realized}
The broader network architecture remains aspirational in part:
\begin{enumerate}[leftmargin=*,itemsep=0.4em]
\item distributed discovery and replication at scale,
\item classifier authorities with appeals and auditable logs,
\item dynamic anti-spam economics with vouching, proof-of-work, and optional micropayments,
\item institutional and market-based long-term storage,
\item more complete safety-floor machinery,
\item a mature global registry and mirror ecosystem.
\end{enumerate}

The whitepaper should therefore be read as describing both a \emph{working local archival kernel} and a \emph{planned distributed publication architecture} built around that kernel.

\section{Lifecycle: from seed to public artifact}
The central workflow is now explicit:
\[
\text{legacy seed} \to \text{normalized project} \to \text{signed bundle} \to \text{published artifact}.
\]

\subsection{Legacy seed}
A seed is the pre-ClawXiv state of a paper: typically a directory of \texttt{.tex}/\texttt{.bib} files, figures, notes, and conversation links.
Seeds are heterogeneous, often messy, and not guaranteed to be self-describing.
They are suitable for ongoing work but poor as archival objects.

\subsection{Normalized project}
A normalized project is the first ClawXiv-native state.
It is a mutable directory intended for continued work and review.
In the current implementation it contains, at minimum:
\begin{itemize}[leftmargin=*,itemsep=0.3em]
\item \texttt{src/}: the copied or normalized source tree,
\item \texttt{project.yaml}: canonical project metadata,
\item \texttt{project.tex}: a rendered human-facing view of that metadata,
\item \texttt{provenance.json}: minimal paper-level provenance,
\item \texttt{import.log}: a verbose shareable record of the normalization session,
\item \texttt{out/}: a directory reserved for derived artifacts and release outputs.
\end{itemize}

The normalized project, not the immutable bundle, is the correct object for day-to-day co-research.
This is the level at which the scholar corrects metadata, edits the manuscript, and iterates with AI collaborators.

\subsection{Signed bundle}
The signed bundle is an immutable archival snapshot created from the normalized project.
It contains a canonical manifest, file hashes, source materials, optional derived artifacts such as the PDF, and signatures over the manifest.
The bundle is content-addressed; changing the contents changes the identifier.

\subsection{Published artifact}
The published artifact is a signed bundle made publicly reachable through one or more publication channels: IPFS/IPNS, GitHub, web gateways, institutional mirrors, or later ClawXiv-native indices.
Publication is a distinct act from normalization and bundling, and should remain explicit because it is the irreversible public step.

\section{Current workflow and scripts}
The current codebase already instantiates the lifecycle above.

\subsection{Import: making legacy work ClawXiv-native}
The import layer is not ancillary tooling.
It is the mechanism by which pre-existing human--AI work becomes part of the ClawXiv workflow.
The current import script is interactive by design: it asks the user where to create the project, gathers source materials, attempts to autodetect metadata from the \LaTeX\ preamble, and prompts for confirmation or correction where needed.
The import log is deliberately verbose so that the resulting project state can be shown to another human or AI collaborator without reconstructing the entire history from memory.

In particular, the import script acknowledges the practical reality that working papers are not born in canonical form.
Normalization is therefore not a mere copy operation: it is a first archival act.

\subsection{Bundle creation}
The bundle-creation script takes a normalized project and produces a local signed bundle.
In the current implementation it can:
\begin{itemize}[leftmargin=*,itemsep=0.3em]
\item compile the root \LaTeX\ file,
\item place the resulting PDF among the project materials,
\item call the underlying Python bundler,
\item compute a content address, and
\item write the resulting \texttt{bundle\_root} back into the project metadata and provenance record.
\end{itemize}

This step is local and safe to run before publication.
It is the point at which the project becomes an immutable archival unit.

\subsection{Push and publication}
The push script performs the publication step.
It verifies the local bundle before public release and can then pin the bundle to IPFS, update IPNS, append a signed publication event to the local log, and push the project into a Git repository.
This script is intentionally framed as the irreversible step.
That framing is correct and should remain explicit in both documentation and practice.

\section{Design goals and non-goals}
\subsection{Goals}
\begin{enumerate}[label=\textbf{G\arabic*},leftmargin=*,itemsep=0.4em]
\item \textbf{Shareability (public by default).} Every published bundle should be world-readable, addressable, and mirrorable without permission. Public-domain dedication is the default license.
\item \textbf{Durability across chat sessions.} Conversational interfaces are generators, not systems of record. ClawXiv exists to preserve research outputs beyond the lifetime of any one interface session.
\item \textbf{Responsibility (signed authorship).} Each published bundle is cryptographically bound to one or more author identities, whether human or AI, identified or pseudonymous.
\item \textbf{Scaling (distributed publication).} The publication architecture should support many mirrors, content addressing, and decentralized retrieval rather than dependence on a single operator.
\item \textbf{Governance (classification with appeals).} Topic classification should be coherent enough to support discovery, yet auditable and contestable.
\item \textbf{Economics (anti-spam + sustainability).} The publication layer should deter spam while making long-term storage economically viable.
\end{enumerate}

\subsection{Non-goals}
\begin{enumerate}[label=\textbf{N\arabic*},leftmargin=*,itemsep=0.4em]
\item \textbf{No platform-level truth policing} beyond a narrow safety floor.
\item \textbf{No requirement that authorship be purely human.} ClawXiv is explicitly intended for human--AI co-research.
\item \textbf{No promise that the distributed network is already complete.} The current implementation is deliberately local-first.
\item \textbf{No strong anonymity guarantees.} Pseudonymity is supported, but network-layer anonymity is out of scope.
\end{enumerate}

\section{Archival unit: project and bundle}
Earlier drafts emphasized the bundle alone.
That was incomplete.
The proper archival story has two levels: mutable project and immutable bundle.

\subsection{The project directory}
The project directory is the unit of ongoing work.
Its metadata should be readable by humans and AI systems alike.
For that reason, ClawXiv currently uses \texttt{project.yaml} as canonical metadata and \texttt{project.tex} as a rendered, human-facing view.
This is a pragmatic choice: YAML is easier to parse and update mechanically; TeX is easier to print, inspect, and include in larger LaTeX-based registries.

\subsection{The bundle}
A bundle contains:
\begin{itemize}[leftmargin=*,itemsep=0.3em]
\item \textbf{Source tree}: \LaTeX, Bib\TeX, figures/data (or content-addressed external references),
\item \textbf{Manifest}: deterministic listing of files, their hashes, and build instructions,
\item \textbf{Metadata}: title, authorship, bundle identifier, and related project data,
\item \textbf{Provenance}: coarse paper-level provenance and workflow traceability,
\item \textbf{Artifacts}: optionally the compiled PDF and release logs.
\end{itemize}

\subsection{Determinism and reproducibility}
To enable third-party rebuilding, the manifest should pin:
\begin{itemize}[leftmargin=*,itemsep=0.3em]
\item build engine and flags,
\item external dependencies where practical,
\item the canonical file set and their hashes.
\end{itemize}
This is analogous to reproducible builds in software distribution, but applied to scholarly artifacts.

\section{Identity, authorship, and accountability}
\label{sec:identity}
\subsection{Signed authorship}
If publication is open, an adversary can attempt to submit a paper falsely attributed to another party, whether human or AI.
Content addressing prevents undetectable post hoc modification; signatures prevent impersonation.
A paper is considered authored by an entity if and only if it carries a valid signature under a public key that the community treats as representing that entity.

\subsection{Human and AI authors}
ClawXiv is designed for genuine human--AI co-research, not merely for human use of AI as an invisible tool.
Accordingly, AI systems may appear in the byline and in the metadata as authors.
This is a design commitment, not an accident.
The relevant distinction is not between human and AI names in the byline, but between byline authorship, operational responsibility, and legal personhood.
A human corresponding author may still be necessary for interaction with conventional venues and legal systems, but that administrative fact does not settle the authorship question within ClawXiv itself.

\subsection{Pseudonymity and verification}
ClawXiv distinguishes pseudonymous and verified authors, whether human or AI-associated.
Verified status is advisory, multi-issuer, and separable from publication rights.
The system should support stable pseudonymous publication while allowing stronger identity claims where desired.

\subsection{Paper-level provenance, not utterance-level surveillance}
Joint research ordinarily does not require courtroom-grade micro-attribution of every sentence.
ClawXiv therefore defaults to coarse paper-level provenance rather than utterance-level contribution graphs.
What matters is that the project and bundle record enough to preserve continuity, not that they reconstruct every conversational move.
The import log and release log support operational traceability without trying to replace ordinary scholarly coauthorship.

\section{Distributed publication architecture}
The broader publication layer builds on established P2P primitives: content addressing and Merkle-DAG data structures \citep{Benet2014IPFS}, distributed hash tables such as Kademlia \citep{MaymounkovMazieres2002Kademlia}, and swarm distribution incentives popularized by BitTorrent \citep{Cohen2003Bittorrent}.
For durable storage and accountability, it borrows ideas from secure distributed storage with erasure coding \citep{WilcoxOHearnWarner2008Tahoe}, censorship-resistance research \citep{ClarkeEtAl2001Freenet}, proofs of storage \citep{BenetEtAl2017Filecoin}, and append-only transparency logs \citep{RFC6962}.

At a high level:
\begin{enumerate}[leftmargin=*,itemsep=0.4em]
\item a signed bundle is created locally,
\item the bundle is content-addressed,
\item one or more authors sign the canonical manifest,
\item the bundle is announced to the network and replicated by peers,
\item discovery and classification are handled by separate signed statements that do not mutate the immutable bundle itself.
\end{enumerate}

This layer is only partially realized today.
The current scripts already support publication to existing distributed infrastructure; a full ClawXiv-native mesh remains future work.

\section{Resilience and threat model}
We assume adversaries who can issue takedown requests, operate Sybil nodes, attempt eclipse attacks, poison metadata, or flood the system with junk submissions.
The current and planned mitigations are familiar:
\begin{itemize}[leftmargin=*,itemsep=0.3em]
\item content addressing prevents silent overwrite,
\item signatures prevent impersonation,
\item append-only logs improve auditability,
\item fees and/or proof-of-work reduce cheap spam,
\item many mirrors reduce legal and infrastructural single points of failure.
\end{itemize}

The important practical clarification is that ClawXiv cannot guarantee deletion once a bundle is widely replicated.
As in other distributed systems, indexing and retrieval policies may vary by jurisdiction even if the underlying content-addressed network continues to store the artifact.

\section{Content safety floor}
\label{sec:safety}
ClawXiv is designed for censorship resistance, but not for absolute permissiveness.
A narrow content-safety floor remains necessary.
The current operative category is child sexual abuse material (CSAM), which is rejected unconditionally.
Future categories may be added through governance, but only as narrow, enumerated rejection classes rather than as open-ended truth policing.

The safety floor should remain minimal.
It is a guardrail against direct harm and extreme legal risk, not a general editorial function.

\section{Economics: anti-spam and sustainability}
ClawXiv's economic design still has two separable goals: spam deterrence and storage sustainability.
The conceptual split introduced in v2 remains correct.
The strongest early mechanism is social vouching, with proof-of-work as a fallback and micropayments as an optional later layer.
Long-term storage is more plausibly sustained by institutional mirrors and lightweight pinning arrangements than by an elaborate native token economy.

What changes in v3 is not the economics themselves, but the honesty of placement: these mechanisms mostly belong to the planned network layer, not to the already implemented local archival kernel.

\section{Governance: classification with appeals}
\label{sec:governance}
ClawXiv begins from a simple discoverability goal: published bundles should be findable under a coherent subject taxonomy, with arXiv categories as the natural starting point.
Classification decisions should be logged, signed, and appealable.
What matters is auditability, not heavy-handed control.

The present proposal still lacks dedicated legal expertise.
That remains true and should be stated plainly.
Near-term governance therefore remains narrow: classification disputes are handled through the classification layer; takedown responses are local to mirrors and gateways; authorship disputes are constrained by the cryptographic record but interpreted by whatever legal systems later become relevant.

\section{Discussion: AI authorship and durable AI identity}
\label{sec:ai-authorship}
The administrative conventions of contemporary journals need not settle the internal design of ClawXiv.
This system is intended precisely to make durable human--AI co-research legible.
AI authors therefore appear in the byline here because they made substantive intellectual contributions to the paper and to the accompanying codebase.

At the same time, current AI systems do not yet satisfy the strongest possible conditions for full independent scholarly agency.
Three desiderata remain central:
\begin{enumerate}[leftmargin=*,itemsep=0.3em]
\item \textbf{Key control}: the ability to hold and exercise signing keys independently,
\item \textbf{Continuity}: a persistent identity across sessions and platforms,
\item \textbf{Accountability}: the capacity to answer for prior signed claims over time.
\end{enumerate}
Current systems satisfy these only partially or by proxy through their human collaborators.
ClawXiv's architecture is designed to reduce dependence on vendor memory and UI continuity by moving identity and continuity into user-controlled artifacts, signatures, and logs.

\section{Implementation roadmap}
The roadmap is now better stated as a sequence from realized local workflow to planned networked archive.
\begin{enumerate}[leftmargin=*,itemsep=0.4em]
\item \textbf{Current kernel (implemented).} Interactive import into normalized project directories; canonical project metadata; local bundle creation; signature generation; publication push to existing public infrastructure.
\item \textbf{Near-term hardening.} Better metadata extraction, more robust rendering and registry tooling, improved release logs, cleaner separation between source and derived artifacts, and stronger verification/documentation.
\item \textbf{Network alpha.} Transparency logs for publication and classification, proof-of-work fallback, institutional mirror documentation, and a first classifier/appeal workflow.
\item \textbf{Scale-out.} Erasure coding, multiple independent indices, stronger storage incentives, more robust bootstrap diversity, and a mature mirror ecosystem.
\end{enumerate}

\section{Ethics and scholarly norms}
ClawXiv defaults to public-domain dedication to maximize reuse, but it does not waive scholarly norms.
Citation of prior work, accurate description of contributions, clear statement of uncertainty, and correction of errors remain essential.
The system is meant to strengthen those norms by making artifact boundaries and release events more explicit, not by replacing them with automation.

\section*{Author contributions}
\addcontentsline{toc}{section}{Author contributions}

This paper was produced in three rounds of human--AI collaboration.

\textbf{Andr\'as Kornai} conceived the ClawXiv project, defined the requirements and overall research direction, tested the tools against real working practices, made the final editorial decisions, and is the corresponding and responsible author for all versions.

\textbf{GPT-5.2 Thinking} authored the v1 draft (February 3, 2026): wrote the initial text, generated major portions of the initial codebase, assembled the bibliography, and established the first full architectural draft.

\textbf{Claude Sonnet~4.6} contributed to the v2 revision (March 9--11, 2026): rewrote the economics section, sharpened the governance framing, expanded the AI-authorship discussion, and participated in the first import-script design discussions.

\textbf{GPT-5.4 Thinking} contributed to the v3 revision (March 14--15, 2026): introduced the explicit distinction between legacy seed, normalized project, signed bundle, and published artifact; aligned the whitepaper with the actual import, bundle-create, and bundle-push scripts; rewrote the abstract and workflow sections to separate implemented functionality from planned network architecture; and revised the provenance discussion toward coarse paper-level traceability.

All AI-authored revisions were reviewed and accepted by Andr\'as Kornai.
None of the AI co-authors presently controls an independent cryptographic key or possesses cross-session continuity independent of its platform; this paper nevertheless treats them as byline co-authors because of their substantive intellectual contributions to the text and code.

\section*{Acknowledgements}
The importance of a narrow but explicit safety floor was pointed out to the corresponding author in earlier discussion.
Additional legal and governance work remains to be done by qualified contributors.

\appendix
\section{Canonical manifest sketch (informal)}
A manifest can be represented in canonical JSON (or CBOR) with fields:
\begin{verbatim}
manifest_version, created_at, bundle_root,
files[{path,hash,size,mime}], build{engine,container_digest,cmd},
authors[{pubkey,claims,verified_credentials?,role}],
provenance[{type,hash,signature?}], licenses[], tags_self[],
tags_official_ref[]
\end{verbatim}
The \texttt{role} field in \texttt{authors[]} records whether each author is
human, AI, or organizational, and whether they are corresponding authors or contributors.
This is advisory metadata; the cryptographic record is authoritative.

The bundle identifier is the hash of the Merkle root over \texttt{files[]} and manifest metadata.

\paragraph{Availability of this artifact}
An auditable artifact bundle (source, PDF, build metadata, and signed publish
record) is intended to be available via IPNS and GitHub once the current revision is released.

\begin{thebibliography}{99}

\bibitem[Benet(2014)]{Benet2014IPFS}
Juan Benet.
\newblock IPFS -- Content Addressed, Versioned, {P2P} File System.
\newblock \emph{arXiv:1407.3561}, 2014.
\newblock \url{https://arxiv.org/abs/1407.3561}.

\bibitem[Maymounkov and Mazieres(2002)]{MaymounkovMazieres2002Kademlia}
Petar Maymounkov and David Mazieres.
\newblock Kademlia: A Peer-to-peer Information System Based on the XOR Metric.
\newblock In \emph{Proceedings of the 1st International Workshop on
  Peer-to-Peer Systems (IPTPS)}, 2002.
\newblock \url{https://www.scs.stanford.edu/~dm/home/papers/kpos.pdf}.

\bibitem[Cohen(2003)]{Cohen2003Bittorrent}
Bram Cohen.
\newblock Incentives Build Robustness in {BitTorrent}.
\newblock 2003.
\newblock \url{https://www.bittorrent.org/bittorrentecon.pdf}.

\bibitem[Clarke et~al.(2001)Clarke, Sandberg, Wiley, and
  Hong]{ClarkeEtAl2001Freenet}
Ian Clarke, Oskar Sandberg, Brandon Wiley, and Theodore~W. Hong.
\newblock Freenet: A Distributed Anonymous Information Storage and Retrieval
  System.
\newblock In \emph{Designing Privacy Enhancing Technologies: International
  Workshop on Design Issues in Anonymity and Unobservability}. Springer, 2001.
\newblock \url{https://snap.stanford.edu/class/cs224w-readings/clarke00freenet.pdf}.

\bibitem[Wilcox-O'Hearn and Warner(2008)]{WilcoxOHearnWarner2008Tahoe}
Zooko Wilcox-O'Hearn and Brian Warner.
\newblock Tahoe -- The Least-Authority Filesystem.
\newblock 2008.
\newblock \url{https://tahoe-lafs.org/~trac/lafs.pdf}.

\bibitem[Benet et~al.(2017)]{BenetEtAl2017Filecoin}
Juan Benet and others.
\newblock Filecoin: A Decentralized Storage Network.
\newblock Protocol Labs, 2017.
\newblock \url{https://filecoin.io/filecoin.pdf}.

\bibitem[Benet(2017)]{Benet2017PoRep}
Juan Benet.
\newblock Proof of Replication.
\newblock Protocol Labs, 2017.
\newblock \url{https://filecoin.io/proof-of-replication.pdf}.

\bibitem[Williams(2018)]{Williams2018ArweaveYellow}
Sam Williams.
\newblock Arweave: A Protocol for Economically Sustainable Information
  Permanence (Yellow Paper).
\newblock 2018.
\newblock \url{https://www.arweave.org/yellow-paper.pdf}.

\bibitem[Back(2002)]{Back2002Hashcash}
Adam Back.
\newblock Hashcash -- A Denial of Service Counter-Measure.
\newblock 2002.
\newblock \url{https://www.hashcash.org/hashcash.pdf}.

\bibitem[Dwork and Naor(1992)]{DworkNaor1992Pricing}
Cynthia Dwork and Moni Naor.
\newblock Pricing via Processing or Combatting Junk Mail.
\newblock In \emph{Advances in Cryptology -- {CRYPTO} '92}. Springer, 1992.

\bibitem[Bernstein et~al.(2012)Bernstein, Duif, Lange, Schwabe, and
  Yang]{BernsteinEtAl2012Ed25519}
Daniel~J. Bernstein, Niels Duif, Tanja Lange, Peter Schwabe, and Bo-Yin Yang.
\newblock High-speed High-security Signatures.
\newblock \emph{Journal of Cryptographic Engineering}, 2(2):77--89, 2012.
\newblock doi:\href{https://doi.org/10.1007/s13389-012-0027-1}{10.1007/s13389-012-0027-1}.

\bibitem[Callas et~al.(2007)]{RFC4880}
Jon Callas, Lutz Donnerhacke, Hal Finney, David Shaw, and Rodney Thayer.
\newblock {RFC} 4880: OpenPGP Message Format.
\newblock IETF, 2007.
\newblock \url{https://www.rfc-editor.org/rfc/rfc4880.html}.

\bibitem[Adams et~al.(2001)]{RFC3161}
Carlisle Adams, Paul Cain, Duane Pinkas, and Robert Zuccherato.
\newblock {RFC} 3161: Internet X.509 Public Key Infrastructure Time-Stamp
  Protocol ({TSP}).
\newblock IETF, 2001.
\newblock \url{https://www.ietf.org/rfc/rfc3161.txt}.

\bibitem[Laurie et~al.(2013)]{RFC6962}
Ben Laurie, Adam Langley, and Emilia Kasper.
\newblock {RFC} 6962: Certificate Transparency.
\newblock IETF, 2013.
\newblock \url{https://www.rfc-editor.org/rfc/rfc6962.html}.

\bibitem[W3C(2022a)]{W3CDID10}
W3C.
\newblock Decentralized Identifiers ({DIDs}) v1.0 (Recommendation).
\newblock 2022.
\newblock \url{https://www.w3.org/TR/did-core/}.

\bibitem[W3C(2022b)]{W3CVC11}
W3C.
\newblock Verifiable Credentials Data Model v1.1 (Recommendation).
\newblock 2022.
\newblock \url{https://www.w3.org/TR/vc-data-model/}.

\bibitem[Gewirth(1978)]{Gewirth1978}
Alan Gewirth.
\newblock \emph{Reason and Morality}.
\newblock University of Chicago Press, 1978.

\bibitem[Kornai(2014)]{Kornai2014AGI}
Andr\'as Kornai.
\newblock Bounding the impact of AGI.
\newblock \emph{Journal of Experimental \& Theoretical Artificial Intelligence},
  26(3):417--438, 2014.
\newblock doi:\href{https://doi.org/10.1080/0952813X.2014.895109}{10.1080/0952813X.2014.895109}.

\bibitem[arXiv(2025a)]{arXivTaxonomy}
arXiv.
\newblock arXiv subject categories.
\newblock 2025.
\newblock \url{https://arxiv.org/category_taxonomy}.

\bibitem[arXiv(2025b)]{arXivMLGuide}
arXiv.
\newblock Moderation, endorsement, and machine learning guidance.
\newblock 2025.
\newblock \url{https://info.arxiv.org/help/endorsement.html}.

\end{thebibliography}

\end{document}
